% BHU PYQ -- Professional Exam Template
\documentclass[11pt,a4paper]{article}

\usepackage[a4paper,top=2.5cm,bottom=2.5cm,left=2.8cm,right=2.8cm,headheight=14pt]{geometry}
\usepackage{amsmath,amssymb}
\usepackage[T1]{fontenc}
\usepackage[utf8]{inputenc}
\usepackage{lmodern}
\usepackage{microtype}
\usepackage{fancyhdr}
\usepackage{enumitem}
\usepackage{hyperref}
\usepackage{lastpage}
\usepackage{parskip}
\usepackage{tikz}
\usetikzlibrary{arrows.meta,shapes,positioning}

\hypersetup{colorlinks=true,urlcolor=black,linkcolor=black,
  pdftitle={B.Sc. (Hons.) Zoology Sem-I ZOB-301 (Old Course 2009-10) -- BHU PYQ}}

\pagestyle{fancy}
\fancyhf{}
\renewcommand{\headrulewidth}{0.4pt}
\renewcommand{\footrulewidth}{0.4pt}
\fancyhead[L]{\small\textit{Banaras Hindu University}}
\fancyhead[R]{\small\textit{Zoology Paper: ZOB-301 (Old)}}
\fancyfoot[C]{\small Page \thepage\ of \pageref{LastPage}}

\newlist{parts}{enumerate}{3}
\setlist[parts,1]{label=(\alph*),leftmargin=*,itemsep=4pt,topsep=3pt}
\setlist[parts,2]{label=(\roman*),leftmargin=*,itemsep=2pt,topsep=2pt}
\setlist[parts,3]{label=(\arabic*),leftmargin=*,itemsep=2pt,topsep=2pt}
\newcommand{\pts}[1]{\hfill\textit{[#1]}}

\begin{document}

\begin{center}
  {\large\bfseries BANARAS HINDU UNIVERSITY}\\[2pt]
  {\normalsize B.Sc. (Hons.) Semester I Examination, 2016-17}\\[6pt]
  \rule{0.8\linewidth}{0.5pt}\\[6pt]
  {\large\bfseries Zoology}\\[4pt]
  {\large\bfseries Paper: ZOB-301 --- Comparative Physiology, Endocrinology and Developmental Biology}\\[2pt]
  {\normalsize\textbf{(Old Course 2009-10 Syllabus)}}\\[4pt]
  {\normalsize Time: Three hours \hfill Full Marks: 70}
\end{center}

\rule{\linewidth}{0.4pt}

\smallskip
\noindent\textit{Note: Answer six questions including Question No. 1 of Section-A, Question No. 1 of Section-B, and Question No. 1 of Section-C which are compulsory. Use a separate answer book for each Section. The figures in the right-hand margin indicate marks.}

\bigskip
\begin{center}
  \textbf{SECTION -- A} \\
  \textbf{(Comparative Physiology) (Marks: 26)}
\end{center}

\smallskip
\noindent\textit{Note: Answer two questions including Question No. 1 which is compulsory.}

\medskip\noindent\textbf{Question 1.} Answer the following parts:
\begin{parts}
  \item Define the following:\pts{$2 \times 1 = 2$}
  \begin{parts}
    \item Ammonotelism
    \item Resting membrane potential.
  \end{parts}

  \item State whether the following statements are True or False with proper reasons:\pts{$3 \times 2 = 6$}
  \begin{parts}
    \item Counter current flow mechanism is found only in fishes.
    \item Shivering is associated with protection against hot environment.
    \item Na$^+$-K$^+$ pump provides an additional contribution to the resting potential.
  \end{parts}
\end{parts}

\medskip\noindent\textbf{Question 2.} Answer the following parts:\pts{$9 + 9 = 18$}
\begin{parts}
  \item Describe the functions of various types of leukocytes found in the blood.
  \item Describe the chemical digestion of proteins in animals.
\end{parts}

\medskip\noindent\textbf{Question 3.} Answer the following parts:\pts{$6 + 6 + 6 = 18$}
\begin{parts}
  \item Discuss the mechanism of branchial respiration in fishes.
  \item Give an account of various kinds of proteins involved in muscle contraction.
  \item Describe ornithine cycle of urea formation.
\end{parts}

\pagebreak
\begin{center}
  \textbf{SECTION -- B} \\
  \textbf{(Endocrinology) (Marks: 26)}
\end{center}

\smallskip
\noindent\textit{Note: Answer three questions including Question No. 1 which is compulsory.}

\medskip\noindent\textbf{Question 1.} Answer the following parts:
\begin{parts}
  \item Write whether the following statements are True or False. Justify your answers:\pts{$4 \times 1 = 4$}
  \begin{parts}
    \item Zona glomerulosa cells of adrenal gland secrete glucocorticoids.
    \item Testosterone is produced by Sertoli cells of the testis.
    \item Glucagon is a hyperglycemic hormone.
    \item Calcitonin is produced by parafollicular cells of the parathyroid gland.
  \end{parts}

  \item Define the following:\pts{$4 \times 1 = 4$}
  \begin{parts}
    \item Endocrine mode of hormone delivery
    \item Zona fasciculata
    \item Graafian follicle
    \item Neurohypophysis.
  \end{parts}
\end{parts}

\medskip\noindent\textbf{Question 2.} Draw a well labelled diagram of Graafian follicle. Write important biological actions of estrogen.\pts{$4 + 5 = 9$}

\medskip\noindent\textbf{Question 3.} Describe the structure of the parathyroid gland. Comment upon the role of parathormone in Ca$^{2+}$ homeostasis.\pts{$4 + 5 = 9$}

\medskip\noindent\textbf{Question 4.} Describe the structure and different parts of pituitary gland and name the hormones secreted by them.\pts{9}

\pagebreak
\begin{center}
  \textbf{SECTION -- C} \\
  \textbf{(Developmental Biology) (Marks: 18)}
\end{center}

\smallskip
\noindent\textit{Note: Answer two questions including Question No. 1 which is compulsory.}

\medskip\noindent\textbf{Question 1.} Answer the following parts:
\begin{parts}
  \item Define the following terms:\pts{$4 \times 0.5 = 2$}
  \begin{parts}
    \item Acrosome reaction
    \item Invagination
    \item Amnion
    \item Differentiation.
  \end{parts}

  \item State whether the following statements are True or False. Justify your answer in 2--3 sentences:\pts{$2 \times 2 = 4$}
  \begin{parts}
    \item Nieuwkoop centre is also called as primary organiser of amphibian.
    \item The function of the Hensen's node is equivalent to the dorsal lip of the amphibian blastopore.
  \end{parts}
\end{parts}

\medskip\noindent\textbf{Question 2.} Answer the following parts:\pts{$6 + 6 = 12$}
\begin{parts}
  \item Define the role of organiser in amphibian embryo development.
  \item Define the process of regeneration and describe the limb regeneration in amphibians.
\end{parts}

\medskip\noindent\textbf{Question 3.} Write short notes on any three of the following:\pts{$3 \times 4 = 12$}
\begin{parts}
  \item External fertilization
  \item Fate map of chick
  \item Fast block to polyspermy
  \item Primitive streak.
\end{parts}

\vfill
\rule{\linewidth}{0.4pt}
\begin{center}\small
\href{https://bhu-exam.vercel.app/}{Click Here}\\[4pt]\href{https://me-aryan.vercel.app/}{Click Here}\\[4pt]\href{https://quashan-me.vercel.app/}{Click Here}
\end{center}

\end{document}
